\documentclass[11pt,a4paper]{article}
\usepackage[margin=2.3cm]{geometry}
\usepackage{graphicx}
\usepackage{booktabs}
\usepackage{amsmath,amssymb}
\usepackage{microtype}
\usepackage[hidelinks]{hyperref}
\usepackage{caption}
\captionsetup{font=small,labelfont=bf}
\setlength{\parskip}{0.4em}
\setlength{\parindent}{0pt}

\title{\textbf{Pre-Sputnik transients on POSS-I plates:\\a reanalysis of the Earth-shadow, nuclear-test, and alignment claims}}
\author{Donal Hill\\\small Personal note, not for publication}
\date{October 2026}

\begin{document}
\maketitle

\begin{abstract}
\noindent
The VASCO project reports short-lived, point-like transients on 1950s POSS-I red plates and three headline statistical
results: a deficit of transients inside Earth's shadow at geostationary altitude (quoted at $\sim$7.6$\sigma$), an excess
of transients around US nuclear tests, and aligned groups of transients on single plates (best candidate $\sim$3.9$\sigma$).
Using the public 107,875-transient catalogue and plate scans, we find that none of the three survives once the
analysis counts plates and nights rather than individual transients, uses correct plate dates, and compares against nulls
that keep the real clustering. The shadow deficit is reproduced as a per-transient binomial ($z=-7.5$) but rests on 9 of
561 plates; plate-level tests give $p\approx0.2$, and POSS-I cannot exclude a vanishing fraction of up to $\sim$0.5, so the
claim is unsupported rather than refuted. The nuclear association is broad in time, appearing as strongly before shots as
after, and disappears once season and galactic latitude are controlled. Aligned groups occur at or below the rate a
clustering-preserving null predicts, and the published selection cut is nearly blind to east--west lines. Along the way,
we find that about three-quarters of POSS-I red-plate dates in the STScI DSS archive are one day early, in contiguous
blocks of the plate serial numbers, verified with asteroid trails, Palomar's night list, and the handwritten plate labels.
\end{abstract}

\section{Background and data}

The claims examined are those of Villarroel et al.\ (2025, PASP; shadow and alignments), Bruehl \& Villarroel (2025,
Sci.\ Rep.; nuclear tests), and Bruehl, Doherty, Streblyanska, \& Villarroel (2026; machine-learning follow-up). The
interpretation offered is that some transients are sunlit reflections from objects in high orbit before any satellites.

\textbf{Data.} The catalogue $V$ of 107,875 transient candidates (Solano et al.\ 2022 pipeline), the plate-based
control sample (635 POSS-I red ``XE'' plates $\times$ 100 random positions with plate UT), and the VASCO analysis code are
public (\texttt{github.com/dca-doherty/VASCO-ML} and \texttt{VASCO-challenge}). Plate headers and scans come from
the STScI DSS FITS archive, and shot times from DOE/NV-209 Rev.~16. Every transient maps unambiguously onto one plate taken
on the same UT date (100\% assignment; 0.14\% in same-night plate overlaps). The ``strict'' sample used in the published
analyses has 99,933 transients on 561 plates and 276 nights.

\textbf{Prior work.} Watters et al.\ (2026, PASA) showed on the smaller public subsets that Poisson errors are
invalid because transient density varies strongly between plates, that $\sim$11 plates touch the shadow, and that
pasting the shadow footprints onto random plates makes a $-8\sigma$ ``deficit'' common. They asked for the 107k sample
to be re-tested this way; that is done here. VASCO's reply (Villarroel et al.\ 2026) predates their Monte Carlo test.
Doherty (2026) independently counted 9 shadow plates, and Hambly \& Blair (2024) argued that the transients are
emulsion defects introduced when the survey plates were copied.

\section{The plate dates in the DSS archive are wrong}

All the timing-sensitive analyses depend on plate dates. Palomar labelled each plate with the night's double date
(``Nov 29/30''), and every POSS-I red exposure (02--12~UT) falls on the \emph{second} date. In long runs of the serial
numbers, the DSS header copied the first date instead.

\begin{figure}[h]
\centering
\includegraphics[width=0.5\linewidth]{figs/label_xe356.png}\\[0.4em]
\includegraphics[width=0.92\linewidth]{figs/asteroids.png}
\caption{Top: the label projected onto plate XE356 at Palomar (``E-441 Nov 29/30 1951''); its DSS header gives 1951-11-29
with a UT clock time of 07:05, which fell on 30 November. Below: known asteroids on DSS cutouts with JPL ephemerides
overlaid (red; crosses hourly). At the header date the predicted tracks are empty (top, XE356 and XE617); one day later,
bright trails sit exactly on the tracks with the length of the exposure (bottom). Control plate XE422 (right column)
shows the opposite.}
\label{fig:evidence}
\end{figure}

Four independent lines of evidence (Figs.~\ref{fig:evidence} and \ref{fig:dates}):
\begin{itemize}\itemsep0pt
\item \textbf{Asteroid trails} on 62 plates (SkyBoT, JPL Horizons from Palomar, DSS cutouts recalibrated against Gaia to
0.2--0.5$''$). Trails appear only at the corrected date, with zero clean detections at impossible offsets ($\pm$12~h,
$\pm$2~d). The header clock time is close to mid-exposure (median offset +4~min).
\item \textbf{Palomar's observing-night list} (VizieR VI/25, VASCO's own ``strict'' file, which uses evening dates)
agrees with the date rule on 320 of 321 decidable plates.
\item \textbf{Handwritten labels} read from the scan margins on 55 plates, the 15 decisive ones checked by eye.
\item \textbf{Serial order}: exposure numbers increase monotonically in time once dates are corrected.
\end{itemize}

\begin{figure}[h]
\centering
\includegraphics[width=0.95\linewidth]{figs/dates.pdf}
\caption{Header-date offset against plate serial number, by evidence type. Shaded: the two blocks where headers are one
day early (E-30 to E-637 and E-968 to E-1610). XE113 and XE389 are the two exceptions to the block rule.}
\label{fig:dates}
\end{figure}

Of the 635 analysis plates, 469 carry a date one day early, 161 are correct, and five are wrong by months or years (XE217,
XE128, XE603, XE498, XE098; fixed from their labels). VASCO's \texttt{obs\_date} inherits these errors, as does any
analysis using DSS epochs. A one-day error moves the anti-solar point by $\sim$1$^\circ$, so it barely affects the
shadow test, but it shifts every night-level timing analysis by a day on three-quarters of plates.

\section{Earth's shadow}

\textbf{Where 7.6$\sigma$ comes from.} 495 in-shadow transients against an expected 0.692\% of 99,933 gives $z=-7.50$
in a one-sided per-transient binomial, the test in VASCO's released code. It treats every transient as an independent
trial. Per-plate counts have variance/mean $\approx300$ (median 131, maximum 4,166), and only 9 of 561 plates touch the
shadow. With corrected dates, 342 transients fall in shadow against 521 expected from plate area: the pooled deficit
comes almost entirely from three plates that simply have few transients overall (Fig.~\ref{fig:shadow}a).

\begin{figure}[h]
\centering
\includegraphics[width=\linewidth]{figs/shadow.pdf}
\caption{(a) Observed minus expected in-shadow transients per shadow plate, if each held the survey-average count. XE617,
XE369, and XE422 carry the deficit; their sunlit parts are sparse too. (b) The in-footprint count when the 9 real shadow
footprints are pasted onto random plates (Watters-style null on $V$).}
\label{fig:shadow}
\end{figure}

\begin{table}[h]
\centering\small
\begin{tabular}{lr}
\toprule
Test (geocentric shadow at GEO, corrected dates) & Result \\
\midrule
Per-transient binomial (published construction) & $z=-7.5$ \\
Footprint-paste null on $V$ (all plates / sky-matched / topocentric) & $p=0.21$ / 0.16 / 0.31 \\
Bootstrap over observing nights & $z\approx-1.8$ \\
Negative binomial across plates (with $|b|$ covariate) & $-2\Delta\ln L=0.16$ \\
Within-plate vanishing fraction $\varphi$ (density ratio, straddling plates) & $0.00$, null s.d.\ 0.34 \\
\bottomrule
\end{tabular}
\caption{The shadow deficit with plates as the unit of analysis.}
\label{tab:shadow}
\end{table}

\textbf{Within plates.} On the 8 plates the shadow edge crosses, the density of transients inside and outside the
shadowed part gives a vanishing fraction $\varphi=0.00$. The estimator is unbiased under injected thinning, but with so
few plates its null spread is 0.34 (geocentric) to 0.45 (topocentric), so the power against $\varphi=0.3$ is only about
50\%. Weighting the expectation for crowding (Section~6) moves $\varphi$ by at most $-0.1$. The data therefore show
no significant deficit, but cannot exclude $\varphi\lesssim0.5$; VASCO's implied $\sim$1/3 is within about $1\sigma$.

\textbf{Geometry.} The Methods of Bruehl et al.\ (2026) describe a topocentric, precessed 8.96$^\circ$ penumbra. The
repository documentation and \texttt{reproduce.py} run a geocentric 8.50$^\circ$ cone at 08:00~UT without precession,
and the published 0.692\% and top-decile 0.31\% reproduce only with that version. Applied consistently, the topocentric
geometry still gives a pooled per-transient deficit (ratio 0.55--0.63), so the mismatch is a documentation problem,
not the cause of the deficit.

\section{Nuclear tests}

The published test compares transient counts on nights within $\pm1$~day of US test dates with all other nights.
With DOE/NV-209 times, Nevada shots fired at $\sim$14~UT and Pacific shots at $\sim$18~UT, so a plate on the same UT date
was exposed hours \emph{before} the shot. With corrected dates, the excess sits on the nights just before shots (1.8$\times$)
and just after (2.1$\times$), at similar levels. The pre-specified $\pm1$-day window gives $p=0.06$--0.10 depending on the
null, and the excess is carried by one observing run (Operation Teapot, March 1955; $p>0.15$ without it). VASCO's test
list also omits all 11 Upshot--Knothole shots (1953) and Teapot MET.

\textbf{Look-elsewhere.} We defined 56 plausible analyses: single-day lags from $-7$ to $+7$, symmetric and one-sided
windows, deliberately arbitrary windows (``4--6 days after''), each for all US tests and for Nevada only. The null is the
real test calendar slid by every whole day from 30 to 1,100 (2,142 calendars). Twenty of the 56 reach local $p<0.05$; the
strongest is ``6 days after a Nevada shot'' (4.0$\times$, $p=0.005$), which no mechanism predicts. The best-of-56 result has
global $p=0.11$, but the \emph{breadth} (20 windows lighting up at once) is unusual ($p=0.007$). Soviet and British test
calendars give 1 and 0 of 28 windows respectively.

\begin{figure}[h]
\centering
\includegraphics[width=\linewidth]{figs/nuclear.pdf}
\caption{(a) Mean transients per observing night by month, with the number of Nevada shots in the survey period.
(b) Mean galactic latitude of each night's plates. (c) How many of the 56 analyses reach local $p<0.05$ for slid test
calendars: sliding by any number of days (grey) breaks the season, sliding by whole years (orange) keeps it.}
\label{fig:nuclear}
\end{figure}

\textbf{Season.} A broad association that ignores timing points to a shared cause through the calendar. Nevada tests
clustered in spring; spring midnight skies lie towards the galactic pole (mean plate $|b|$ 46--58$^\circ$ in March--May
against $\sim$20$^\circ$ from July to December); and transient counts per night track galactic latitude (Spearman 0.44;
Fig.~\ref{fig:nuclear}). Holding the season fixed by sliding the calendar only by whole years ($\pm$20 days) leaves 7 of
56 analyses at $p<0.05$ (breadth $p=0.15$, global $p=0.26$), and adjusting counts for galactic latitude leaves 2--3
(global $p\approx0.2$). Watters et al.\ speculated about seasonality through the observing schedule; the mechanism here
is the sky, and it removes the effect.

\section{Aligned transients}

VASCO searched the 298k parent sample in boxes of a few to 30$'$ and kept groups whose RA and Dec have Pearson
$|r|>0.99$, then estimated each candidate's chance probability for that single box with the uniform-density formula of
Edmunds \& George (1985). The paper notes that the survey-wide probability cannot be estimated.

\begin{figure}[h]
\centering
\includegraphics[width=\linewidth]{figs/alignments.pdf}
\caption{(a) For every pair of transients within 10$'$ on a plate, the perpendicular offset of each third transient lying
between them. Pure clustering gives a smooth distribution near zero offset; real alignments give a spike (orange: the
same with 100 synthetic aligned groups injected). (b) The fraction of perfect 4-point lines passing VASCO's
$|r(\mathrm{RA},\mathrm{Dec})|>0.99$ cut, by orientation.}
\label{fig:align}
\end{figure}

\begin{itemize}\itemsep0pt
\item \textbf{Orientation bias.} The Pearson cut passes a perfect east--west line 1\% of the time (24\% at 2$^\circ$,
72\% at 5$^\circ$, $\sim$99\% for diagonals). Glints from an equatorial geostationary object drifting against the stars
would line up east--west.
\item \textbf{No excess collinearity} in $V$: $+208\pm146$ configurations within 3$''$ of a line against a linear
baseline, where injecting 100 aligned groups gives a 3.3$\sigma$ spike (Fig.~\ref{fig:align}a). Near-perfect lines have
the same orientation distribution as the control.
\item \textbf{Counts at the chance rate.} Against a null that jitters every transient by up to 1$'$ (keeping the clustering,
destroying arcsecond collinearity), the survey holds 4,399 four-point alignments within a 5$''$ strip (null 4,744) and 1,931
five-point alignments within 10$''$ (null 2,238). Configurations like the shortlisted candidates occur thousands of times.
\end{itemize}

\section{What the catalogue is made of}

\textbf{Plate edges.} Transient density rises from the plate centre (1.4$\times$ at 1.5--2$^\circ$, 2.0--2.6$\times$ at
2--2.75$^\circ$) while the control footprint stays flat, in 88\% of plates. A sky-fixed population cannot do this, so a large
plate-fixed component is present (emulsion, copying, scanning, or cross-match failures where images degrade).

\begin{figure}[h]
\centering
\includegraphics[width=\linewidth]{figs/crowding.pdf}
\caption{(a) Transients per plate against galactic latitude. (b) At the same low latitude in northern fields (no
airmass difference), the dense inner Galaxy yields far fewer transients than the sparse anticentre.}
\label{fig:crowding}
\end{figure}

\textbf{Crowding.} Counts per plate rise about 2.2\% per degree of $|b|$ ($\sim$3.8$\times$ from the plane to
$|b|=60^\circ$), and at the same latitude the dense inner Galaxy (Aquila--Cygnus) yields 3.6$\times$ fewer transients
than the sparse anticentre (medians 24.5 and 89.5 per plate; Fig.~\ref{fig:crowding}). A candidate survives only if no
modern catalogue source lies within 5$''$, and chance matches are far more likely in crowded fields. A physical flash
population would not know where the Milky Way is; stellar flares would do the opposite. This selection is what links
test seasons to transient counts.

None of this establishes what any individual transient is. Busko (2026) reports transients sharing the optical coma of
neighbouring stars, which points to light through the telescope for at least some objects. The decisive test would be
to compare the DSS and SuperCOSMOS scans, made from different glass copies, for all 107k candidates.

\section{Conclusions}

\begin{itemize}\itemsep0pt
\item \textbf{Shadow:} the 7.6$\sigma$ is a per-transient counting artefact. With plates as the unit, $p\approx0.2$; within
plates, $\varphi=0.00$ but POSS-I cannot exclude values up to $\sim$0.5. Unsupported, not refuted.
\item \textbf{Nuclear:} the association is broad in time, as strong before shots as after, and vanishes once season and
galactic latitude are controlled.
\item \textbf{Alignments:} no more than chance plus clustering produces; the selection cut cannot see east--west lines.
\item \textbf{Archive:} about three-quarters of POSS-I red-plate dates in the DSS archive are one day early, in serial blocks,
with a handful of larger errors; this is independent of the transient debate and worth reporting to STScI.
\end{itemize}

\textbf{Caveats.} The top machine-learning decile (a 55\% deficit) cannot be tested because the public scores are
degenerate; its scores include plate-level features, so it shares the clustering problem, but we do not claim it is wrong.
The PASP significance was reproduced from the released code rather than the paper's own method description. We read the
arXiv version of Bruehl et al.\ (2026), not the journal text. Small-sample $p$-values (6--9 plates, a few observing
runs) are approximate. An adversarial internal review of every claim is in \texttt{redteam/report.md}.

\textbf{Acknowledgements.} The Digitized Sky Surveys were produced at the Space Telescope Science Institute under U.S.
Government grant NAG W-2166, based on photographic data obtained using the Oschin Schmidt Telescope on Palomar Mountain.
This work used the VASCO catalogue and code released by B.\ Doherty and collaborators (CC~BY~4.0 / MIT), the Spanish
Virtual Observatory, SkyBoT, JPL Horizons, and Gaia DR3 via VizieR.

\small
\begin{thebibliography}{99}\itemsep0pt
\bibitem{bv25} Bruehl, S., \& Villarroel, B. 2025, Sci.\ Rep., doi:10.1038/s41598-025-21620-3
\bibitem{bdsv26} Bruehl, S., Doherty, B., Streblyanska, A., \& Villarroel, B. 2026, Sci.\ Rep.\ 16, 31890; arXiv:2604.18799
\bibitem{busko} Busko, I. 2026, arXiv:2606.08319
\bibitem{doe} DOE/NV-209 Rev.~16 2015, \emph{United States Nuclear Tests, July 1945 through September 1992}
\bibitem{doherty} Doherty, B. 2026, arXiv:2604.00056
\bibitem{eg85} Edmunds, M.~G., \& George, G.~H. 1985, MNRAS 213, 905
\bibitem{hb24} Hambly, N.~C., \& Blair, A. 2024, RAS Techniques and Instruments
\bibitem{solano} Solano, E., Villarroel, B., \& Rodrigo, C. 2022, MNRAS 515, 1380
\bibitem{v25} Villarroel, B., et al.\ 2025, PASP 137, 104504
\bibitem{v26} Villarroel, B., Streblyanska, A., Bruehl, S., \& Geier, S. 2026, arXiv:2602.15171
\bibitem{watters} Watters, W.~A., et al.\ 2026, PASA, doi:10.1017/pasa.2026.10230; arXiv:2601.21946
\end{thebibliography}

\end{document}
